\section{Mesure initiale : affichage vs calcul}

Avant toute parallélisation, il faut distinguer le coût du rendu graphique du coût du calcul physique. Le visualiseur affiche déjà deux temps :
\begin{itemize}
    \item \texttt{Render time}
    \item \texttt{Update time}
\end{itemize}

Cette séparation est importante, car seule la partie calcul est directement concernée par l'étape Numba. Si l'affichage reste significatif, la loi d'Amdahl limitera le speed-up global observé dans la simulation interactive.

\subsection{Définition exacte des temps mesurés}

Les deux temps affichés dans la boucle interactive ne correspondent pas à un profilage fin au sens d'un outil dédié ; ils représentent une découpe simple de la boucle principale du visualiseur.

Dans \texttt{stages/01\_numba\_threads/visualizer3d.py}, trois instants sont relevés à chaque itération :
\begin{itemize}
    \item \(t_1\) : fin de l'itération précédente ;
    \item \(t_3\) : juste après l'appel à \texttt{\_render()} ;
    \item \(t_2\) : juste après \texttt{update\_points(updater(dt))}.
\end{itemize}

On obtient alors :
\[
\texttt{Render time} = t_3 - t_1
\qquad\text{et}\qquad
\texttt{Update time} = t_2 - t_3
\]

Autrement dit :
\begin{itemize}
    \item \texttt{Render time} inclut la gestion des événements SDL du début de boucle, le rendu OpenGL du frame courant et l'échange de buffers ;
    \item \texttt{Update time} inclut l'appel au calcul physique \texttt{updater(dt)} ainsi que la copie des nouvelles positions dans le visualiseur via \texttt{update\_points} ;
    \item \texttt{Update time} n'inclut ni la gestion d'événements ni le rendu ;
    \item la première itération est fortement perturbée par l'initialisation et, dans le cas Numba, par la compilation JIT.
\end{itemize}

En conséquence, ces deux temps donnent une bonne vision macroscopique de la boucle interactive, mais ils ne constituent pas un profilage détaillé instruction par instruction.

\subsection{Mesures observées}

Les premières mesures ont été relevées sur la version de référence \texttt{nbodies\_grid\_numba.py}, avec \texttt{data/galaxy\_1000}, \(dt = 0.0015\) et une grille \(15 \times 15 \times 1\). Après l'itération initiale de compilation JIT, les valeurs stables observées sont d'environ \SI{5}{ms} pour le rendu et \SI{52}{ms} pour la mise à jour physique.

\begin{table}[H]
    \centering
    \begin{tabular}{l c}
        \toprule
        \textbf{Mesure initiale} & \textbf{Valeur} \\
        \midrule
        Temps de rendu & \(\approx\) \SI{5}{ms} \\
        Temps de mise à jour & \(\approx\) \SI{52}{ms} \\
        Partie dominante & calcul des trajectoires \\
        \bottomrule
    \end{tabular}
    \caption{Mesure initiale pour comparer affichage et calcul}
    \label{tab:initial_render_update}
\end{table}

\begin{figure}[H]
    \centering
    \begin{tikzpicture}
\begin{axis}[
    ybar,
    bar width=22pt,
    width=0.82\linewidth,
    height=6.2cm,
    ymin=0,
    ylabel={Temps (ms)},
    symbolic x coords={Render + events, Update},
    xtick=data,
    xticklabel style={font=\small, align=center},
    nodes near coords,
    every node near coord/.append style={font=\footnotesize},
    enlarge x limits=0.35,
    grid=both,
    minor y tick num=1,
    axis line style={draw=black!60},
]
\addplot[fill=enstaBleuClair, draw=enstaBleuFonce] coordinates {
    (Render + events, 5.0)
    (Update, 52.0)
};
\end{axis}
\end{tikzpicture}

    \caption{Comparaison des deux temps relevés dans la boucle interactive de référence}
    \label{fig:initial_render_update}
\end{figure}
